\section{Capital: What You Can Actually Afford to Lose}

\subsection{Distinguishing 4 Types of Capital}

Founders tend to think of capital as money in a bank account or the amount an investor might give them. At the beginning, a more useful question is: \textbf{How much money can I responsibly commit to an experiment without threatening obligations I cannot abandon?}

A person with ₹10 lakh in cash and ₹9 lakh of near-term obligations may have less deployable capital than someone with ₹2 lakh and stable income, low expenses, and no urgent debt. A person with a laptop, development skills, and 8 hours a week may test a software service without cash outlay. Someone starting food manufacturing cannot assume the same.

You need to distinguish 4 economic categories:

\begin{itemize}[leftmargin=11pt,itemsep=1.5pt,topsep=1.5pt,parsep=0pt]
  \item \textbf{Wealth:} Estimated value of assets minus obligations. Not necessarily available to spend.
  \item \textbf{Liquidity:} Cash or assets convertible to cash promptly without severe haircut.
  \item \textbf{Investable cash:} Money remaining after strictly protected reserves and known debts.
  \item \textbf{Affordable loss:} The smaller sum you consciously accept putting at risk on 1 experiment.
\end{itemize}

Do not treat personal borrowing capacity, a residential plot, or a future salary as equivalent to free startup cash. \textbf{Cap the downside before excitement alters judgment.}

\subsection{The Personal Financial Boundary}

First, construct a cash-flow baseline independent of the new venture:

\begin{tcolorbox}[
  enhanced,
  colback=white,
  colframe=BayesBorder!80!white,
  boxrule=0.4pt,
  arc=2pt,
  left=4pt,right=4pt,top=4pt,bottom=4pt,
  before skip=3pt,after skip=4pt
]
\scriptsize
\setlength{\tabcolsep}{3pt}
\renewcommand{\arraystretch}{1.15}
\begin{tabularx}{\linewidth}{@{} X r @{}}
\toprule
\textbf{Monthly Cash-Flow Line Item} & \textbf{Amount} \\
\midrule
Net reliable personal income & ₹\rule{2.2cm}{0.3pt} \\
Essential household expenses & (₹\rule{2.0cm}{0.3pt}) \\
Existing debt servicing and EMIs & (₹\rule{2.0cm}{0.3pt}) \\
Insurance, dependents, medical reserve & (₹\rule{2.0cm}{0.3pt}) \\
Known near-term contractual provisions & (₹\rule{2.0cm}{0.3pt}) \\
\midrule
\textbf{Monthly surplus before venture spend} & \textbf{₹\rule{2.0cm}{0.3pt}} \\
\bottomrule
\end{tabularx}
\end{tcolorbox}

Second, audit liquid reserves to find provisional deployable cash:

\begin{tcolorbox}[
  enhanced,
  colback=white,
  colframe=BayesBorder!80!white,
  boxrule=0.4pt,
  arc=2pt,
  left=4pt,right=4pt,top=4pt,bottom=4pt,
  before skip=3pt,after skip=5pt
]
\scriptsize
\setlength{\tabcolsep}{3pt}
\renewcommand{\arraystretch}{1.15}
\begin{tabularx}{\linewidth}{@{} X r @{}}
\toprule
\textbf{Liquid Resource Audit} & \textbf{Amount} \\
\midrule
Unrestricted personal liquid cash & ₹\rule{2.2cm}{0.3pt} \\
Protected 6-month emergency reserve & (₹\rule{2.0cm}{0.3pt}) \\
Near-term taxes and statutory obligations & (₹\rule{2.0cm}{0.3pt}) \\
Expected short-term household shortfalls & (₹\rule{2.0cm}{0.3pt}) \\
\midrule
\textbf{Provisional Deployable Capital} & \textbf{₹\rule{2.0cm}{0.3pt}} \\
\bottomrule
\end{tabularx}
\end{tcolorbox}

\subsection{Decide the Loss Before Forecasting Return}

Suppose a founder has ₹1.8 lakh of genuinely discretionary cash after reserves. They might decide that the 1st experiment cannot put more than ₹18,000 at risk, even if an imagined opportunity could earn crores. The point is to cap the downside before optimism overrides discipline.

An affordable-loss budget includes more than direct purchases:

\begin{equation}
\begin{aligned}
&\text{Affordable Loss} \\
&\quad = \text{Cash Spend} \\
&\qquad + \text{Irreversible Commitments} \\
&\qquad + \text{Opportunity Cost of Time} \\
&\qquad + \text{Downside Liabilities}.
\end{aligned}
\end{equation}

For each experiment, specify:
\begin{itemize}[leftmargin=11pt,itemsep=1.5pt,topsep=1.5pt,parsep=0pt]
  \item Maximum cash you will spend.
  \item Maximum founder hours you will commit.
  \item A clear calendar decision date.
  \item Evidence required to continue versus conditions that require stopping.
  \item Obligations that remain if the experiment produces 0 revenue.
\end{itemize}

\subsection{Capital Beyond Cash}

Ventures can launch with customer pre-orders, customer-financed pilots, leased gear, or sweat equity. Each has trade-offs: customer advances create strict delivery debts; leasing introduces fixed overhead; pre-sales create refund risk. None makes economic cost vanish.

\textbf{Seek the cheapest honest experiment that could falsify the business hypothesis: you are buying truth, not the appearance of progress.}
